\documentclass[11pt]{article}
\usepackage[T1]{fontenc}
\usepackage[utf8]{inputenc}
\usepackage{lmodern}
\usepackage[margin=1in]{geometry}
\usepackage{amsmath,amssymb,amsthm,mathtools}
\usepackage{microtype,booktabs,tabularx,array,enumitem}
\usepackage[round,authoryear]{natbib}
\usepackage{xcolor,hyperref}
\definecolor{linkblue}{HTML}{174B71}
\hypersetup{colorlinks=true,linkcolor=linkblue,citecolor=linkblue,
  urlcolor=linkblue,pdftitle={Optimal Randomized Complexity of Fixed Non-Euclidean Proximal Oracles},
  pdfauthor={Colombo Team}}
\urlstyle{same}
\setlength{\emergencystretch}{2em}
\setlength{\parskip}{0pt}
\setlength{\parindent}{1.25em}
\setlist[enumerate]{leftmargin=*,itemsep=3pt,topsep=5pt}
\setlist[itemize]{leftmargin=*,itemsep=3pt,topsep=5pt}
\newtheorem{theorem}{Theorem}[section]
\newtheorem{lemma}[theorem]{Lemma}
\newtheorem{proposition}[theorem]{Proposition}
\newtheorem{corollary}[theorem]{Corollary}
\theoremstyle{definition}
\newtheorem{definition}[theorem]{Definition}
\theoremstyle{remark}
\newtheorem{remark}[theorem]{Remark}
\newcommand{\R}{\mathbb R}
\newcommand{\E}{\mathbb E}
\newcommand{\PP}{\mathbb P}
\newcommand{\eps}{\varepsilon}
\newcommand{\norm}[1]{\left\lVert#1\right\rVert}
\newcommand{\ip}[2]{\left\langle#1,#2\right\rangle}
\newcommand{\F}{\mathcal F}
\newcommand{\D}{\mathcal D}
\newcommand{\Qrand}{Q_{\mathrm{rand}}}
\newcommand{\Qdet}{Q_{\mathrm{det}}}
\DeclareMathOperator*{\argmin}{arg\,min}
\DeclareMathOperator{\supp}{supp}
\newcolumntype{Y}{>{\raggedright\arraybackslash}X}
\title{Optimal Randomized Complexity of\\
Fixed Non-Euclidean Proximal Oracles}
\author{Colombo Team\\[3pt]
  \texttt{colombo.2026@outlook.com}}
\date{}

\begin{document}
\maketitle

\begin{abstract}
We study convex minimization using an exact proximal oracle that returns
$\argmin_x\{f(x)+\lambda\norm{x-c}_p^s\}$ at a queried center $c$, where
$2<p<s<\infty$ and $\lambda>0$ are fixed. The objective is finite and convex
and has a minimizer of $p$-norm at most $R$. We prove a lower bound for
arbitrary randomized algorithms, allowing unbounded query centers and any
measurable final output. For every query budget $T$, a finite distribution
of objectives forces error
$\Omega_s(\lambda R^s/(T+1)^{s+s/p-1})$ with probability at least $11/12$;
dimension $O_p((T+1)^{p+3})$ suffices. The proof combines randomly hidden
coordinates with a common radial term and an exact truncation identity for
proximal responses. Together with the deterministic upper bound of Adil,
Bullins, Jambulapati, and Sidford, this establishes
$\Theta_{p,s}(1+(\lambda R^s/\eps)^{1/(s+s/p-1)})$ query complexity in
sufficiently large dimension. Thus the known rate remains optimal without
support or span restrictions and without a logarithmic loss. We also verify
that the upper bound uses only the fixed-coefficient proximal replies,
including for nonsmooth objectives.
\end{abstract}

\noindent\textbf{Keywords:} oracle complexity; convex optimization;
non-Euclidean proximal oracles; randomized lower bounds; acceleration.
\par\medskip

\begin{samepage}
\noindent\textbf{AI Usage.}
Nearly the entire research pipeline for this paper was carried out by
Colombo, powered by GPT-5.6 Sol.
\par
\end{samepage}

\section{Introduction}\label{sec:introduction}

A proximal oracle solves a regularized optimization problem at each query.
The central complexity question is how many such solves are needed to
minimize the original objective. We consider the non-Euclidean oracle
\begin{equation}\label{eq:intro-oracle}
P_f(c)=\argmin_{x\in\R^d}
\bigl\{f(x)+\lambda\norm{x-c}_p^s\bigr\},
\qquad 2<p<s<\infty,
\end{equation}
where the coefficient $\lambda>0$ is fixed throughout the interaction.
The algorithm chooses the center $c$ and receives only the minimizer.
The unknown objective is finite and convex, and at least one of its
minimizers lies in the $p$-norm ball of radius $R$ about the origin.

For this oracle, \citet[Theorem~4]{ABJS} give a deterministic algorithm
whose query complexity is
\begin{equation}\label{eq:intro-rate}
O_{p,s}\!\left(1+\kappa^{1/\rho}\right),
\qquad
\kappa=\frac{\lambda R^s}{\eps},
\qquad
\rho=s+\frac{s}{p}-1.
\end{equation}
Their lower bound establishes the same accuracy exponent for algorithms
whose iterates satisfy a support restriction. A general oracle algorithm,
however, can probe any coordinate, choose centers of arbitrarily large
norm, and output any measurable function of its transcript and random seed.
Whether these additional freedoms improve the rate is a separate question.

We show that the rate in \eqref{eq:intro-rate} remains optimal for arbitrary
randomized algorithms in sufficiently large dimension. The lower bound
uses a finite distribution of functions fixed before the algorithm's
random seed is drawn. It therefore applies to adaptive randomized queries
in the usual black-box model, with no restriction on the support, span,
or norm of the query centers. A separate argument covers arbitrary final
outputs. The dimension needed for a $T$-query lower bound is
$O_p((T+1)^{p+3})$; all constants and quantifiers are given in
Section~\ref{sec:model}.

\subsection{Comparison with prior work}\label{sec:related}

Table~\ref{tab:comparison} compares the results for the same exact
proximal oracle and the common parameter range $2<p<s<\infty$.
The known upper rate and the support-restricted lower rate are due to
\citet{ABJS}; our contribution is the lower bound for unrestricted
randomized algorithms. The upper bound holds in every dimension, and
our lower bound is proved under an explicit dimension condition.

\begin{table}[t]
\centering
\small
\setlength{\tabcolsep}{5pt}
\renewcommand{\arraystretch}{1.18}
\begin{tabularx}{\linewidth}{@{}p{0.22\linewidth}Yp{0.24\linewidth}@{}}
\toprule
Result & Algorithmic scope & Query bound\\
\midrule
\citet{ABJS}\newline Theorem~4
& Deterministic upper bound, in every dimension
& $O_{p,s}(1+\kappa^{1/\rho})$\\
\citet{ABJS}\newline Theorem~5
& Lower bound for proximal zero-respecting iterates
& $\Omega_{p,s}(\kappa^{1/\rho})$\\
Theorems~\ref{thm:budget} and~\ref{thm:matching}
& Lower bound for arbitrary randomized algorithms and arbitrary outputs
& $\Omega_{p,s}(1+\kappa^{1/\rho})$\\
\bottomrule
\end{tabularx}
\caption{Comparison for the fixed oracle \eqref{eq:intro-oracle}.
The lower bounds are high-dimensional statements. Our sufficient condition
is $d\ge\lceil C_p(1+\kappa^{(p+3)/\rho})\rceil$.
The support restriction in the second row is stated explicitly in
\eqref{eq:prior-support}.}
\label{tab:comparison}
\end{table}

The formal restriction in \citet[Definition~22]{ABJS} is
\begin{equation}\label{eq:prior-support}
\supp(c_t)\subseteq\bigcup_{j<t}\supp(P_f(c_j)).
\end{equation}
Requiring a query to lie in the span of earlier replies implies
\eqref{eq:prior-support}, but the support condition permits a larger
algorithm class. We use the latter condition when describing the prior
result. In contrast, our lower bound imposes neither condition and
does not require the output to be an iterate.

The construction builds on the shifted maximum of affine functions and
the radial extension used by \citet[Definition~24 and Eq.~(4)]{ABJS}.
Randomized hiding arguments for convex optimization in general norms
were developed by \citet{DG} for local oracles. Their framework supplies
important context, but a local-oracle lower bound does not directly
control \eqref{eq:intro-oracle}: its reply is the solution of a global
minimization problem. The exact equality of proximal replies established
in Lemma~\ref{lem:exact} addresses this distinction.

Two other related oracle models help delimit the comparison.
\citet{CJJJLST} study a Euclidean ball optimization oracle, with a query
returning a minimizer over a ball around its center, and obtain matching
bounds up to logarithmic factors.
\citet{CGMR} develop non-Euclidean accelerated methods and lower bounds
for high-order smooth convex optimization, including randomized and
parallel local-oracle lower bounds. These results concern local derivative
information or fixed-radius ball minimization. Neither directly yields a
lower bound for the fixed exact proximal oracle studied here.

\subsection{Proof overview}\label{sec:overview}

The hard objective is the maximum of $T+1$ affine functions placed on
randomly chosen signed coordinates, together with a common radial term.
The affine offsets increase along the sequence. A query can expose the
next affine function, but, with high probability, its remaining hidden
coordinates are too small to expose the later ones. Three features make
this argument valid for the unrestricted exact oracle.

First, the hard objectives are globally Lipschitz. Their proximal replies
therefore remain within a fixed distance of their centers, even when the
centers are unbounded. For centers near the minimizer ball, the increasing
offsets keep all omitted affine functions below a retained one at the
truncated proximal minimizer. For distant centers, the radial term has
the same effect. Global majorization then proves that the full and
truncated objectives have exactly the same proximal minimizer. This
identity holds for the complete real-valued reply vector.

Second, adaptivity is handled through a surrogate transcript defined for
every sampled objective. At time $t$, the surrogate answers using only
the first $t$ hidden coordinates. Each query is consequently measurable
with respect to the previously exposed prefix. Probability bounds are
conditioned on that prefix, and the actual and surrogate transcripts are
coupled afterwards. This order preserves the uniform distribution of the
unseen coordinates.

Third, one hidden coordinate remains after the last query. The radial
term rules out outputs outside the radius-$R$ ball. A counting argument
then shows that an arbitrary output inside the ball has a small component
on the last hidden coordinate with high probability, which forces a
nonzero objective gap.

The resulting exponent follows from three scales. With $m=T+1$, the
offset spacing is $\alpha\asymp R/m^{1+1/p}$. Restricting the proximal
displacement to a constant multiple of $\alpha$ requires a Lipschitz
scale $\beta\asymp_s\lambda\alpha^{s-1}$. The comparator improves the
objective by order $\beta R/m^{1/p}$, giving
\[
\frac{\beta R}{m^{1/p}}
\asymp_s\frac{\lambda R^s}{m^{(s-1)(1+1/p)+1/p}}
=\frac{\lambda R^s}{m^\rho}.
\]
Random coordinate hiding determines the sufficient dimension, while
these scales determine the accuracy exponent.

\paragraph{Organization.}
Section~\ref{sec:model} states the model and main results.
Section~\ref{sec:construction} constructs the hard family and proves
exact proximal locality. Section~\ref{sec:randomization} establishes
the randomized lower bound. Section~\ref{sec:upper} verifies the
compatibility of the known upper bound with our oracle and derives the
matching minimax rate. Section~\ref{sec:discussion} discusses the scope
and the remaining dimension question.

\section{Oracle model and main results}\label{sec:model}

Throughout, $2<p<s<\infty$ and $\lambda,R>0$ are fixed.
We write $p^*=p/(p-1)$ for the dual exponent, $e_j$ for the $j$-th
coordinate vector, and $[d]=\{1,\ldots,d\}$.
The notation $O_{p,s}$, $\Omega_{p,s}$, and $\Theta_{p,s}$ suppresses
positive constants depending only on $p$ and $s$.

\subsection{Functions, queries, and success criterion}

For each integer $d\ge1$, let
\begin{equation}\label{eq:function-class}
\F_{p,R}^{d}
=\bigl\{f:\R^d\to\R:\ f\text{ is convex and }
\argmin f\cap\{x:\norm{x}_p\le R\}\ne\varnothing\bigr\}.
\end{equation}
The functions in this class are finite and continuous; no smoothness or
uniform Lipschitz constant is assumed. Write $f^\star=\min_x f(x)$.
The algorithm knows $p,s,\lambda,R,d$, and the target accuracy $\eps>0$.
Its only access to the unknown objective is through
\begin{equation}\label{eq:oracle}
P_f(c)=\argmin_{x\in\R^d}
\bigl\{f(x)+\lambda\norm{x-c}_p^s\bigr\}.
\end{equation}
This minimizer exists and is unique: $f$ is bounded below, while the
regularizer is coercive and strictly convex. Each call returns the
minimizer itself, without a separate function value or subgradient.
The same coefficient $\lambda$ is used in every call.

A randomized algorithm has an independent random seed $U$. Its query
centers and final output are measurable functions of $U$ and the
previous oracle replies. Centers may lie anywhere in $\R^d$, and the
final output need not be a query or a reply. We charge only oracle
calls, and the query budget is a worst-case bound over instances and
random seeds.

Let $\Qrand(\eps;d)$ be the smallest budget for which an algorithm
returns $\widehat x$ satisfying
\begin{equation}\label{eq:success}
\PP_U\{f(\widehat x)-f^\star\le\eps\}\ge\frac23
\qquad\text{for every }f\in\F_{p,R}^{d}.
\end{equation}
Define $\Qdet(\eps;d)$ analogously, requiring deterministic success on
every instance. We suppress $p,s,\lambda,R$ in this notation and use
$\rho$ and $\kappa$ as defined in \eqref{eq:intro-rate}. In particular,
$\kappa$ is a scale ratio, not a condition number.

\subsection{A distributional lower bound}

For an integer $T\ge0$, define
\begin{equation}\label{eq:dimension-T}
\begin{split}
D_p(T)
&=T+\left\lceil
12\bigl(T(T+1)+1\bigr)64^p(T+1)^{p+1}
\right\rceil,\\
c_s&=\frac{s}{4\cdot64^{s-1}},
\qquad
\Delta_T=c_s\frac{\lambda R^s}{(T+1)^\rho}.
\end{split}
\end{equation}
Thus $D_p(T)=O_p((T+1)^{p+3})$.

\begin{theorem}[Unrestricted randomized lower bound]\label{thm:budget}
Let $T\ge0$ be an integer and $d\ge D_p(T)$. There is a finite
distribution $\D_{T,d}$ supported on $\F_{p,R}^{d}$ such that every
randomized algorithm $\mathcal A$ making at most $T$ queries satisfies
\begin{equation}\label{eq:distributional}
\PP_{f\sim\D_{T,d},\,U}
\bigl\{f(\mathcal A^{P_f}(U))-f^\star>\Delta_T\bigr\}
\ge\frac{11}{12}.
\end{equation}
The distribution is independent of the algorithm and its seed.
Consequently, for each such algorithm there is a fixed
$f\in\F_{p,R}^{d}$ on which its failure probability at accuracy
$\Delta_T$ is at least $11/12$.
\end{theorem}

The distribution is given explicitly in \eqref{eq:hard-family}.
The conclusion concerns a single function drawn before any queries are
made. In particular, it does not use an adversary that changes the
objective during the interaction.

\subsection{Matching query complexity}

\begin{theorem}[High-dimensional minimax complexity]\label{thm:matching}
Let $C_p=2+12\cdot64^p$. If
\begin{equation}\label{eq:dimension-corollary}
d\ge\left\lceil C_p\bigl(1+\kappa^{(p+3)/\rho}\bigr)\right\rceil,
\end{equation}
then
\begin{equation}\label{eq:matching}
\Qrand(\eps;d)=\Theta_{p,s}\!\left(1+\kappa^{1/\rho}\right),
\qquad
\Qdet(\eps;d)=\Theta_{p,s}\!\left(1+\kappa^{1/\rho}\right).
\end{equation}
The deterministic upper bound holds for every $d\ge1$.
\end{theorem}

The dimension condition is sufficient for our construction and is not
claimed to be necessary. The theorem establishes a matching accuracy
dependence in growing dimension; it does not characterize minimax
complexity at each fixed dimension. The upper bound is the acceleration
result of \citet{ABJS}; Section~\ref{sec:upper} explains why its
implementation conforms to \eqref{eq:oracle} for the entire function
class \eqref{eq:function-class}.

\section{A hard family with exact proximal locality}\label{sec:construction}

\subsection{Construction and basic properties}

Fix a query budget $T\ge0$, assume $d\ge D_p(T)$, and let $m=T+1$.
We use the parameters
\begin{equation}\label{eq:parameters}
\alpha=\frac{R}{8m^{1+1/p}},
\qquad
\delta=\frac{\alpha}{8},
\qquad
\beta=s\lambda\delta^{s-1},
\qquad
\theta=\frac{\alpha}{8R}.
\end{equation}
In particular,
\begin{equation}\label{eq:B}
B:=\theta^{-p}=64^p m^{p+1},
\qquad
\theta\le\frac1{64},
\qquad
\delta\le\frac R{64}.
\end{equation}
The offset $\alpha$ separates successive affine functions, and $\delta$
will bound the distance from a query center to its proximal reply.

Choose $(J_1,\ldots,J_m)$ uniformly from the ordered samples of $m$
distinct coordinates in $[d]$. Independently, choose
$\sigma_1,\ldots,\sigma_m$ uniformly from $\{-1,+1\}$. Set
\begin{equation}\label{eq:hard-family}
f_{J,\sigma}(x)
=\beta\max\left\{
\max_{1\le i\le m}\{\sigma_i x_{J_i}-i\alpha\},\,
\norm{x}_p-R
\right\}.
\end{equation}
We call the common term $\norm{x}_p-R$ the \emph{radial wall}.
It belongs to the objective and places no restriction on query centers.
The resulting finite distribution is $\D_{T,d}$.
For $1\le t\le m$, define the prefix truncation
\begin{equation}\label{eq:truncated}
f^{[t]}(x)
=\beta\max\left\{
\max_{1\le i\le t}\{\sigma_i x_{J_i}-i\alpha\},\,
\norm{x}_p-R
\right\}.
\end{equation}
Then $f^{[m]}=f_{J,\sigma}$ and $f^{[t]}\le f^{[m]}$ everywhere.

\begin{lemma}[Membership in the function class and a comparison point]\label{lem:legal}
Every $f_{J,\sigma}$ is finite, convex, coercive, and globally
$\beta$-Lipschitz with respect to $\norm{\cdot}_p$. Every minimizer has $p$-norm
strictly smaller than $R$. Moreover,
\begin{equation}\label{eq:optimal-value}
f_{J,\sigma}^\star\le-\frac{\beta R}{2m^{1/p}}.
\end{equation}
The same Lipschitz property holds for every truncated function $f^{[t]}$.
\end{lemma}

\begin{proof}
Each affine function in the maximum has Lipschitz constant one, as does
$\norm{x}_p-R$. The maximum is convex and one-Lipschitz, proving the
corresponding claims after multiplication by $\beta$.
The inequality $f_{J,\sigma}(x)\ge\beta(\norm{x}_p-R)$ implies coercivity.
Since $\alpha<R$,
\[
f_{J,\sigma}(0)=-\beta\alpha<0,
\qquad
\norm{x}_p\ge R\ \Longrightarrow\ f_{J,\sigma}(x)\ge0.
\]
Continuity and coercivity give attainment, and all minimizers lie strictly
inside the radius-$R$ ball.

Consider
\begin{equation}\label{eq:comparator}
\bar x=-\frac{R}{2m^{1/p}}\sum_{i=1}^m\sigma_i e_{J_i}.
\end{equation}
The coordinates are distinct, so $\norm{\bar x}_p=R/2$. Every affine face at
$\bar x$ is at most $-R/(2m^{1/p})$, while the radial wall is $-R/2$,
which is also at most $-R/(2m^{1/p})$. Evaluating the maximum gives
\eqref{eq:optimal-value}.
\end{proof}

\subsection{Uniform displacement and exact truncation}

\begin{lemma}[Uniform displacement]\label{lem:displacement}
If $h$ is a finite convex $\beta$-Lipschitz function and
$z=\argmin_x\{h(x)+\lambda\norm{x-c}_p^s\}$, then
\begin{equation}\label{eq:displacement}
\norm{z-c}_p\le\left(\frac{\beta}{s\lambda}\right)^{1/(s-1)}.
\end{equation}
In particular, the proximal reply of every $f^{[t]}$ is within distance
$\delta$ of its center, irrespective of the size of that center.
\end{lemma}

\begin{proof}
Every $v\in\partial h(z)$ has $\norm{v}_{p^*}\le\beta$: applying the
subgradient inequality in any direction and the Lipschitz inequality bounds
the dual pairing in both signs by $\beta$ times the direction's norm.
If $z=c$ there is nothing to show. Otherwise Fermat's rule gives
\[
v=-s\lambda\norm{z-c}_p^{s-p}J_p(z-c)\in\partial h(z),
\qquad (J_p(u))_j=|u_j|^{p-2}u_j.
\]
Its dual norm is
\[
s\lambda\norm{z-c}_p^{s-p}\norm{J_p(z-c)}_{p^*}
=s\lambda\norm{z-c}_p^{s-1},
\]
since $\norm{J_p(u)}_{p^*}=\norm{u}_p^{p-1}$. Comparing this with $\beta$
proves \eqref{eq:displacement}.
\end{proof}

\begin{lemma}[Equality after truncation]\label{lem:exact}
Fix $1\le t\le m$ and $c\in\R^d$. Suppose
\begin{equation}\label{eq:coordinate-event}
|c_{J_i}|\le\theta\max\{R,\norm{c}_p\}
\quad\text{for every }i=t,\ldots,m.
\end{equation}
Then the exact proximal replies satisfy
\begin{equation}\label{eq:exact-reply}
P_{f^{[m]}}(c)=P_{f^{[t]}}(c).
\end{equation}
\end{lemma}

\begin{proof}
Let $z=P_{f^{[t]}}(c)$ and $C=\norm{c}_p$. Lemma~\ref{lem:displacement}
gives $\norm{z-c}_p\le\delta$. We show that every omitted affine face lies
below one of the retained terms at $z$.

\emph{Case 1: $C\le2R$.}
The coordinate condition and $\theta=\alpha/(8R)$ imply
$|c_{J_i}|\le\alpha/4$ for all $i\ge t$. Therefore
\[
|z_{J_i}|\le\frac{\alpha}{4}+\delta=\frac{3\alpha}{8}
\quad(i\ge t).
\]
The retained $t$-th affine face obeys
\[
\sigma_tz_{J_t}-t\alpha\ge-t\alpha-\frac{3\alpha}{8},
\]
whereas every omitted face $i\ge t+1$ obeys
\[
\sigma_i z_{J_i}-i\alpha
\le\frac{3\alpha}{8}-(t+1)\alpha
=-t\alpha-\frac{5\alpha}{8}.
\]
The retained face exceeds every omitted face by at least $\alpha/4$.

\emph{Case 2: $C>2R$.}
For every omitted face, the coordinate condition gives
\[
\sigma_i z_{J_i}-i\alpha\le|z_{J_i}|\le\theta C+\delta.
\]
The radial wall, which is retained in every truncation, satisfies
\[
\norm{z}_p-R\ge C-\delta-R.
\]
Its difference from the upper bound on any omitted face is at least
\begin{align*}
(1-\theta)C-R-2\delta
&>(1-2\theta)R-2\delta\\
&\ge\frac{15R}{16}>0,
\end{align*}
where \eqref{eq:B} was used. This case includes centers of arbitrarily
large norm.

In either case $f^{[m]}(z)=f^{[t]}(z)$. Since $f^{[m]}\ge f^{[t]}$ globally,
for every $x$ we have
\begin{align*}
f^{[m]}(x)+\lambda\norm{x-c}_p^s
&\ge f^{[t]}(x)+\lambda\norm{x-c}_p^s\\
&\ge f^{[t]}(z)+\lambda\norm{z-c}_p^s\\
&=f^{[m]}(z)+\lambda\norm{z-c}_p^s.
\end{align*}
Uniqueness of the proximal minimizer proves \eqref{eq:exact-reply}.
\end{proof}

\begin{remark}[Exact responses]
The identity in Lemma~\ref{lem:exact} holds for the full real-valued
response vector. The argument uses equality of the two objectives at the
truncated minimizer and global majorization elsewhere. It therefore
controls the exact global minimization in the oracle without any
finite-precision assumption.
\end{remark}

\section{Adaptive indistinguishability and the lower bound}\label{sec:randomization}

\subsection{A surrogate transcript}

Fix a deterministic algorithm with at most $T$ queries. If it stops early,
retain its output and pad the run with queries at zero until the budget is
exactly $T$. This does not change its output or its guarantee.

Define a \emph{surrogate experiment} on the entire probability space of
$(J,\sigma)$ as follows. At time $t=1,\ldots,T$, let $c_t$ be the query
chosen by the algorithm from the preceding surrogate replies, and answer with
\begin{equation}\label{eq:surrogate}
z_t=P_{f^{[t]}}(c_t).
\end{equation}
Let $y$ be its final output and let
\[
\mathcal H_t=\sigma(J_1,\ldots,J_t,\sigma_1,\ldots,\sigma_t).
\]
The query $c_t$ is $\mathcal H_{t-1}$-measurable, the reply $z_t$ is
$\mathcal H_t$-measurable, and $y$ is $\mathcal H_T$-measurable. The
surrogate functions need not equal the full function away from the points
used in Lemma~\ref{lem:exact}.

\subsection{Uniform control of hidden coordinates}

For any vector $c$, define its set of large coordinates
\[
\mathcal B(c)=\{j:|c_j|>\theta\max\{R,\norm{c}_p\}\}.
\]
A deterministic counting inequality gives
\begin{equation}\label{eq:count}
|\mathcal B(c)|
\le\frac{\norm{c}_p^p}{\theta^p\max\{R,\norm{c}_p\}^p}
\le\theta^{-p}=B.
\end{equation}
This bound is uniform over all norms and all directions of $c$.

Let $E_t$ be the event
\[
E_t=\bigl\{|(c_t)_{J_i}|\le\theta\max\{R,\norm{c_t}_p\}
\text{ for every }i=t,\ldots,m\bigr\}.
\]
Conditional on $\mathcal H_{t-1}$, the remaining coordinates are a uniform
ordered sample without replacement from $d-t+1$ positions. Hence
\begin{equation}\label{eq:stage-prob}
\PP(E_t^c\mid\mathcal H_{t-1})
\le\frac{(m-t+1)B}{d-t+1}.
\end{equation}

The final output needs its own event:
\begin{equation}\label{eq:output-event}
H=\{\norm{y}_p\ge R\}\ \cup\ \{|y_{J_m}|\le\alpha/8\}.
\end{equation}
Conditional on $\mathcal H_T$, either $\norm{y}_p\ge R$ and $H$ is
certain, or $\norm{y}_p<R$. In the latter case at most
\[
\left(\frac{8R}{\alpha}\right)^p=B
\]
coordinates of $y$ have absolute value larger than $\alpha/8$. The final
coordinate $J_m$, with $m=T+1$, is uniform on the remaining $d-T$ positions.
Therefore
\begin{equation}\label{eq:output-prob}
\PP(H^c\mid\mathcal H_T)\le\frac{B}{d-T}.
\end{equation}

Let $G=H\cap\bigcap_{t=1}^T E_t$. Taking expectations in
\eqref{eq:stage-prob} and \eqref{eq:output-prob}, and then applying a union
bound, gives
\begin{align}
\PP(G^c)
&\le B\sum_{t=1}^T\frac{m-t+1}{d-t+1}+\frac{B}{d-T}\notag\\
&\le\frac{(Tm+1)B}{d-T}\le\frac1{12}.
\label{eq:good-prob}
\end{align}
The last inequality is exactly the dimension requirement
\eqref{eq:dimension-T}.

The conditioning in this calculation is essential. Each probability
bound is conditioned only on the revealed coordinate and sign prefix,
not on the earlier events $E_1,\ldots,E_{t-1}$. Those events can reveal
information about the unexposed coordinates. Defining the surrogate
experiment on the whole sample space allows us to establish the
prefix-conditional bounds first and intersect the events afterwards.

\subsection{Coupling the transcripts}

On $\bigcap_{t=1}^T E_t$, the actual transcript obtained from the one fixed
function $f^{[m]}$ equals the surrogate transcript. This follows by induction:
identical earlier replies imply identical next queries, and
Lemma~\ref{lem:exact} then gives identical next replies. Consequently, on $G$
the actual output equals the surrogate output $y$ as well.

\subsection{Error of the final output}

\begin{lemma}[The final hidden coordinate]\label{lem:terminal}
On the event $H$ of \eqref{eq:output-event},
\begin{equation}\label{eq:terminal-gap}
f_{J,\sigma}(y)-f_{J,\sigma}^\star
>\frac{\beta R}{4m^{1/p}}.
\end{equation}
\end{lemma}

\begin{proof}
If $\norm{y}_p\ge R$, the wall gives $f_{J,\sigma}(y)\ge0$, and
\eqref{eq:optimal-value} gives a gap of at least
$\beta R/(2m^{1/p})$.

Otherwise $|y_{J_m}|\le\alpha/8$. The last affine face yields
\[
f_{J,\sigma}(y)
\ge\beta(\sigma_m y_{J_m}-m\alpha)
\ge-\beta(m+1/8)\alpha.
\]
Combining this with \eqref{eq:optimal-value} and the choice of $\alpha$,
\begin{align*}
f_{J,\sigma}(y)-f_{J,\sigma}^\star
&\ge\frac{\beta R}{m^{1/p}}
\left(\frac12-\frac18-\frac1{64m}\right)\\
&\ge\frac{23}{64}\frac{\beta R}{m^{1/p}}
>\frac{\beta R}{4m^{1/p}}.
\end{align*}
No property of $y$ other than its norm and one hidden coordinate is used.
\end{proof}

\subsection{Completing the distributional argument}

\begin{proof}[Proof of Theorem~\ref{thm:budget}]
For a fixed deterministic algorithm, \eqref{eq:good-prob}, equality of the
actual and surrogate transcripts on $G$, and
Lemma~\ref{lem:terminal} show that the probability of error exceeding
$\beta R/(4m^{1/p})$ is at least $11/12$. The parameters satisfy
\begin{align*}
\frac{\beta R}{4m^{1/p}}
&=\frac{s\lambda}{4}
\left(\frac{R}{64m^{1+1/p}}\right)^{s-1}
\frac{R}{m^{1/p}}\\
&=\frac{s}{4\cdot64^{s-1}}
\frac{\lambda R^s}{m^{(s-1)(1+1/p)+1/p}}
=\Delta_T,
\end{align*}
because $(s-1)(1+1/p)+1/p=s+s/p-1=\rho$.

For a randomized algorithm, condition on any fixed seed. All of the preceding
inequalities hold with the same distribution $\D_{T,d}$ and the same constants,
regardless of that seed. Integrating over the independent seed proves
\eqref{eq:distributional}. Since $\D_{T,d}$ is finite, at least one function
in its support has failure probability at least the average $11/12$. The selected instance is fixed for that randomized algorithm and does
not depend on the realized seed.
\end{proof}

\section{The upper bound and matching complexity}\label{sec:upper}

The lower bound is self-contained. To obtain the matching characterization,
we use the acceleration theorem of \citet{ABJS}. This section verifies
the information supplied by our oracle, the use of a known radius bound,
and the constant-query regime.

\subsection{Implementing the known acceleration method}

\begin{proposition}[Upper bound of \citet{ABJS} in the present model]
\label{prop:upper}
For every $d\ge1$, the oracle \eqref{eq:oracle} suffices to achieve
\begin{equation}\label{eq:upper}
\Qdet(\eps;d)
\le C_{p,s}\left(1+\left(\frac{\lambda R^s}{\eps}\right)^{1/\rho}\right)
\end{equation}
for a constant $C_{p,s}>0$ depending only on $p$ and $s$.
\end{proposition}

\begin{proof}[Verification of oracle and radius compatibility]
We use \citet[Theorems~4 and~14, Algorithm~2, and Appendix~C]{ABJS}.
Let $z=P_f(c)$ and $u=z-c$. The proximal optimality condition gives
the explicitly computable subgradient
\begin{equation}\label{eq:recover-g}
g=-s\lambda\norm{u}_p^{s-p}J_p(u)\in\partial f(z),
\qquad
(J_p(u))_j=|u_j|^{p-2}u_j.
\end{equation}
If $u=0$, then $0\in\partial f(z)$ and $z$ is a global minimizer.
Otherwise, set $r=\norm{u}_p$ and $\eta=s\lambda r^{s-p}$. Then
\begin{equation}\label{eq:gms}
\ip{g}{u}=-\eta r^p,
\qquad
\norm{g}_{p^*}=\eta r^{p-1},
\qquad
r=\left(\frac{\eta}{s\lambda}\right)^{1/(s-p)}.
\end{equation}
These identities satisfy the exact generalized Monteiro--Svaiter
conditions used in \citet[Lemma~39]{ABJS}, with movement parameter
$\mu^s=s\lambda$. The algorithm uses the returned point, its associated
parameter $\eta$, and linear terms involving $g$. Each is determined
by the proximal reply. Its potential proof uses the subgradient
inequality for the objective, so \eqref{eq:recover-g} also makes the
argument valid for finite nonsmooth convex functions. Function values
are not used to choose queries or to stop the method. The changing
estimates of $\eta$ in that algorithm leave the coefficient $\lambda$
of every proximal call unchanged.

For the radius bound, normalize the objective as
\begin{equation}\label{eq:normalization}
h(v)=\frac{f(Rv)}{\lambda R^s}.
\end{equation}
It has a minimizer $v^\star$ of norm at most one. Its unit-coefficient
proximal oracle is simulated exactly by
$P_h(u)=R^{-1}P_f(Ru)$. With mirror function
$\psi(v)=\norm{v}_p^p$ and initial point zero, the initial Bregman
potential is $\Psi_0=\norm{v^\star}_p^p\le1$.
Thus the phase analysis in \citet[Appendix~C.3]{ABJS} can use the
known upper bound $1$ in both its parameter schedule and its potential
budget; it does not need the actual minimizer distance.

For clarity, this substitution can be read directly from the phase
bounds. Put $D=ps-p+s=p\rho$ and $\mu^s=s$. During a phase starting
at a positive accumulated weight $A$, take the parameter estimate
\begin{equation}\label{eq:phase-parameter}
\bar\eta=A^{-(s-p)(p+1)/D}\mu^{sp^2/D}.
\end{equation}
The two types of iterations in the cited potential analysis require at
most
\begin{equation}\label{eq:phase-count}
K_p\left(1+(A\bar\eta)^{1/p}
+\frac{\mu^{sp/(s-p)}}{A\bar\eta^{s/(s-p)}}\right)
\le K_{p,s}(1+A^{p/D})
\end{equation}
calls to double the weight. Here the numerator of the last term uses
$\Psi_0\le1$. Both nonconstant terms have exponent $p/D$ in $A$.

The initialization also has a uniform bound. If $z=P_h(0)\ne0$,
convexity and \eqref{eq:recover-g} give
$\ip{J_p(z)}{z-v^\star}\le0$, whence
$\norm{z}_p\le\norm{v^\star}_p\le1$ by H\"older's inequality.
The first oracle parameter is therefore at most $s$.
Algorithm~2 of \citet{ABJS}, initialized with parameter estimate one,
then has first positive weight at least $3^{-p}/s$.
If $z=0$, the objective is already minimized.
Summing \eqref{eq:phase-count} over the doubling phases until the
weight reaches $\kappa$ gives
$O_{p,s}(1+\kappa^{p/D})$ calls. The potential guarantee is
$h(x)-h^\star\le\Psi_0/A\le1/A$, so stopping at $A\ge\kappa$
attains error at most $1/\kappa$ for $h$, or $\eps$ for $f$.
Since $p/D=1/\rho$, this gives \eqref{eq:upper}.
\end{proof}

\subsection{The constant-query regime}

If $\eps\ge\lambda R^s$, one query is sufficient: for any minimizer
$x^\star$ with $\norm{x^\star}_p\le R$,
\begin{equation}\label{eq:one-query}
f(P_f(0))+\lambda\norm{P_f(0)}_p^s
\le f(x^\star)+\lambda\norm{x^\star}_p^s
\le f^\star+\lambda R^s.
\end{equation}
The additive constant in \eqref{eq:matching} is also necessary because
the function class has no uniform scale bound on the objective.

\begin{lemma}[At least one query is necessary]\label{lem:zero}
For every $\eps,\lambda,R>0$ and every $d\ge1$,
$\Qrand(\eps;d)\ge1$.
\end{lemma}

\begin{proof}
Choose $M>\eps/R$ and consider
\[
f_+(x)=M|x_1-R|,
\qquad f_-(x)=M|x_1+R|.
\]
Both are finite convex functions with a minimizer at norm $R$. Their
$\eps$-sublevel sets are disjoint, because the corresponding intervals for
$x_1$ have half-width $\eps/M<R$. A zero-query output has the same
distribution on both instances. Its success probabilities on the two
instances therefore sum to at most one, so one of them has success probability
at most $1/2$.
\end{proof}

Combining Lemma~\ref{lem:zero} with \eqref{eq:one-query} gives the exact
constant-query conclusion
\[
\Qrand(\eps;d)=\Qdet(\eps;d)=1
\qquad\text{if }\eps\ge\lambda R^s,
\]
in every dimension $d\ge1$.

\subsection{From the budget bound to the minimax rate}

\begin{proof}[Proof of Theorem~\ref{thm:matching}]
Proposition~\ref{prop:upper} gives the upper bound. For the lower bound, write
\[
b=\left(\frac{c_s}{2}\right)^{1/\rho},
\qquad A=b\kappa^{1/\rho},
\qquad N=\lfloor A\rfloor.
\]
Here $0<b<1$. If $N\ge2$, take $T=N-1$. Then
\[
\Delta_T=c_s\frac{\lambda R^s}{N^\rho}
\ge c_s\frac{\lambda R^s}{A^\rho}=2\eps.
\]
Theorem~\ref{thm:budget} rules out a $T$-query algorithm with success
probability $2/3$ at accuracy $\eps$. Thus
$\Qrand(\eps;d)\ge N\ge A/2$, provided $d\ge D_p(N-1)$.

This dimension requirement follows from \eqref{eq:dimension-corollary}:
since $N\ge2$ and $N\le A\le\kappa^{1/\rho}$,
\begin{align*}
D_p(N-1)
&\le \left\lceil(1+12\cdot64^p)N^{p+3}\right\rceil\\
&\le \left\lceil(1+12\cdot64^p)\kappa^{(p+3)/\rho}\right\rceil.
\end{align*}

If $N<2$, then $A<2$, and Lemma~\ref{lem:zero} gives
$\Qrand(\eps;d)\ge1>A/2$. In both cases,
\[
\Qrand(\eps;d)
\ge\max\left\{1,\frac b2\kappa^{1/\rho}\right\}
\ge\frac b4\left(1+\kappa^{1/\rho}\right).
\]
Finally $\Qrand\le \Qdet$, so the same lower bound and the
deterministic upper bound prove both statements in \eqref{eq:matching}.
\end{proof}

\section{Discussion}\label{sec:discussion}

Theorem~\ref{thm:matching} shows that unrestricted randomization does
not improve the accuracy exponent of fixed non-Euclidean proximal
optimization in sufficiently large dimension. The conclusion includes
arbitrary query centers and outputs, and it holds for exact responses.
The central step is the truncation identity: although the oracle solves
a global optimization problem, a common radial wall and a uniform
displacement bound make its reply depend only on a prefix of the hidden
affine functions with high probability.

The dimension bound has a transparent origin. Each query has at most
$B=64^p(T+1)^{p+1}$ coordinates above the normalized hiding threshold.
Controlling all query--hidden-coordinate pairs contributes a further
factor of order $(T+1)^2$. This yields the sufficient dimension
$O_p((T+1)^{p+3})$ without changing the objective-error exponent.
Table~\ref{tab:examples} illustrates the resulting scales for several
fixed choices of $p$ and $s$.

\begin{table}[htbp]
\centering
\small
\renewcommand{\arraystretch}{1.15}
\begin{tabular}{@{}ccccc@{}}
\toprule
$p$ & $s$ & $\rho$ & Query complexity & Sufficient dimension\\
\midrule
$3$ & $4$ & $13/3$ & $\Theta(\eps^{-3/13})$ & $O(\eps^{-18/13})$\\
$3$ & $6$ & $7$ & $\Theta(\eps^{-1/7})$ & $O(\eps^{-6/7})$\\
$4$ & $8$ & $9$ & $\Theta(\eps^{-1/9})$ & $O(\eps^{-7/9})$\\
\bottomrule
\end{tabular}
\caption{Small-accuracy consequences for $\lambda=R=1$.
The last column gives a sufficient construction size with an appropriately
large constant; it is not a necessary dimension threshold.}
\label{tab:examples}
\end{table}

Several questions remain outside the present result. The proof does
not determine the smallest dimension supporting the lower bound, or
the minimax complexity when $d$ stays fixed as $\eps$ decreases.
It also treats $p$ and $s$ as fixed and does not optimize constants
as these parameters vary. Finally, the model counts exact proximal
calls; it does not measure the computational cost of implementing
one such call or impose additional smoothness on the objective.
Reducing the dimension needed for unrestricted lower bounds is a
natural next problem.

\begin{thebibliography}{9}

\bibitem[Adil et~al.(2026)Adil, Bullins, Jambulapati, and Sidford]{ABJS}
Deeksha Adil, Brian Bullins, Arun Jambulapati, and Aaron Sidford.
\newblock Convex optimization with $p$-norm oracles.
\newblock In \emph{Proceedings of the 37th International Conference on
Algorithmic Learning Theory}, volume 313 of \emph{Proceedings of Machine
Learning Research}, pages 1--38, 2026.
\newblock \url{https://proceedings.mlr.press/v313/adil26b.html}.

\bibitem[Carmon et~al.(2020)Carmon, Jambulapati, Jiang, Jin, Lee, Sidford, and Tian]{CJJJLST}
Yair Carmon, Arun Jambulapati, Qijia Jiang, Yujia Jin, Yin Tat Lee,
Aaron Sidford, and Kevin Tian.
\newblock Acceleration with a ball optimization oracle.
\newblock In \emph{Advances in Neural Information Processing Systems},
volume 33, pages 19052--19063, 2020.
\newblock \url{https://arxiv.org/abs/2003.08078}.

\bibitem[Contreras et~al.(2025)Contreras, Guzm\'an, and Mart\'inez-Rubio]{CGMR}
Juan Pablo Contreras, Crist\'obal Guzm\'an, and David Mart\'inez-Rubio.
\newblock Non-Euclidean high-order smooth convex optimization.
\newblock Extended abstract in \emph{Proceedings of the 38th Annual
Conference on Learning Theory}, volume 291 of \emph{Proceedings of
Machine Learning Research}, page 1330, 2025.
\newblock Full version: arXiv:2411.08987.
\newblock \url{https://proceedings.mlr.press/v291/contreras25a.html}.
\newblock \url{https://arxiv.org/abs/2411.08987}.

\bibitem[Diakonikolas and Guzm\'an(2020)]{DG}
Jelena Diakonikolas and Crist\'obal Guzm\'an.
\newblock Lower bounds for parallel and randomized convex optimization.
\newblock \emph{Journal of Machine Learning Research}, 21(5):1--31, 2020.
\newblock \url{https://jmlr.org/papers/v21/19-771.html}.

\end{thebibliography}
\end{document}
